\section{Analisi dataflow}

\subsection{CFGs}
Nelle analisi di tipo dataflow i programmi vengono rappresentati come
control flow graph (CFG), dove i nodi rappresentano gli statement del
programma e gli archi rappresentano i possibili flussi di controllo tra di
essi.

\begin{esempio}{CFG di un programma semplice}{cfg_example}

Consideriamo il seguente programma:
\begin{verbatim}
1 y = a - b ;
2 x = y + b ;
3 while ( y > a + b )
4   y = y - x * x;
\end{verbatim}

La sua rappresentazione come CFG è la seguente:
\begin{figure}[H]
  \centering
  \includegraphics[width=0.6\textwidth]{images/th_04/01.png}
  \caption{L'immagine illustra il CFG del programma di esempio.}
  \label{fig:th_04_01}
\end{figure}

\end{esempio}

\begin{definizione}{Struttura di un CFG}{cfg_structure}
Dato un programma $P$, il suo control flow graph (CFG) è formato da:
\begin{itemize}
    \item Un insieme di nodi (blocchi) identificati da $i \in \mathbb{N}$
    \item Un insieme di archi rappresentati come coppie $(i_1, i_2)
    \in \mathbb{N} \times \mathbb{N}$ dove $i_1, i_2$ sono nodi del CFG.
\end{itemize}
Dato un CFG, è possibile definire:
\begin{itemize}
    \item $nodes(G)$ come l'insieme di tutti gli indici dei blocchi in $G$.
    \item $stmt (G, i)$ lo statement associato al blocco $i$ in $G$.
    \item $edges(G) \subseteq \mathbb{N} \times \mathbb{N}$ come l'insieme di
    tutti gli archi in $G$.
\end{itemize}
\end{definizione}

\subsection{Analisi di dataflow}

Le analisi di dataflow sono tecniche tradizionali utilizzate principalmente
per le ottimizzazioni nei compilatori.
Queste analisi si basano sulla rappresentazione del programma tramite
CFG e utilizzano un approccio formale basato su sistemi di equazioni.

Le equazioni vengono costruite definendo le relazioni
tra gli stati entranti (\textit{entry}) e uscenti (\textit{exit}) di ogni
blocco del CFG.

\begin{nota}{Scopo dell'analisi}{scopo_df}
La proprietà di interesse è solitamente semplice e orientata a ottimizzare
l'esecuzione a runtime (ad esempio, evitando di calcolare due volte la stessa
espressione).
\end{nota}

\begin{nota}{Classificazione delle analisi di dataflow}{types_of_df}
Le analisi dataflow classiche possono essere classificate secondo due
criteri principali:
\begin{itemize}
    \item \textbf{Direzione del flusso}: indica se l'analisi procede
    seguendo il flusso di controllo del programma (\textbf{forward}) o
    contro di esso (\textbf{backward}).
    \item \textbf{Semantica della proprietà}: indica se l'analisi
    determina ciò che \textbf{può} accadere (\textbf{may/possible}) o ciò che
    \textbf{deve} accadere (\textbf{must/definite}).
\end{itemize}

Il tipo di classificazione determina importanti aspetti formali:

\begin{itemize}
    \item \textbf{Forward vs Backward}:
    \begin{itemize}
        \item Forward: $entry(i)$ dipende da $exit(j)$ dei predecessori
        \item Backward: $exit(i)$ dipende da $entry(j)$ dei successori
    \end{itemize}
    
    \item \textbf{May vs Must}:
    \begin{itemize}
        \item May (possible): si utilizza l'unione ($\cup$) al join point
        (basta che la proprietà valga lungo \emph{un} cammino)
        \item Must (definite): si utilizza l'intersezione ($\cap$) al join point
        (la proprietà deve valere lungo \emph{tutti} i cammini)
    \end{itemize}
\end{itemize}

Queste caratteristiche determinano la struttura del reticolo e l'operatore
di join utilizzato nell'analisi.

\end{nota}

\subsubsection{Available expressions}

\begin{definizione}{Available expressions}{available_expressions}
``For each program point, which expressions must have already
been computed, and not later modified, on all paths to the program
point''

Formalmente un'espressione $e$ è disponibile in un punto del programma $p$ se
e solo se $e$ è stata calcolata lungo ogni cammino che porta a $p$ e nessuna
delle variabili che occorrono in $e$ è stata ridefinita dopo il calcolo di $e$.
\end{definizione}

\begin{nota}{Obiettivi dell'analisi}{ae_objectives}
Quest'analisi è un'analisi \textbf{forward} e \textbf{must} che ha come
obiettivi principali:
\begin{itemize}
    \item Eliminazione di sottoespressioni comuni (Common Subexpression
    Elimination): evitare di ricalcolare espressioni il cui valore è già
    disponibile.
    \item Allocazione ottimale dei registri: memorizzare le espressioni più
    frequentemente utilizzate in registri temporanei.
\end{itemize}
\end{nota}

\paragraph{Dominio dell'analisi}

Sia $AE$ l'insieme di tutte le espressioni aritmetiche non banali che
compaiono nel programma. Lo stato dell'analisi in ogni punto del programma è
un elemento di $\mathcal{P}(AE)$, ovvero un sottoinsieme delle espressioni
aritmetiche.

\begin{definizione}{Funzioni ausiliarie}{ae_auxiliary_functions}
È possibile definire le seguenti funzioni ausiliarie:
\begin{itemize}
    \item $AExp : \textit{Expr} \to \mathcal{P}(AE)$ restituisce l'insieme
    delle espressioni aritmetiche non banali contenute in un'espressione $e$.
    \item $AExpB : \textit{BExpr} \to \mathcal{P}(AE)$ restituisce l'insieme
    delle espressioni aritmetiche non banali contenute in un'espressione
    booleana $b$.
    \item $FV : AE \to \mathcal{P}(\textit{Var})$ restituisce l'insieme delle
    variabili libere (free variables) che occorrono in un'espressione $a$.
\end{itemize}
\end{definizione}

\paragraph{Funzioni di trasferimento}

Per ogni statement del programma, è necessario definire le funzioni $gen$ e
$kill$ che descrivono come lo statement modifica l'insieme delle espressioni
disponibili.

\begin{definizione}{Gen e Kill}{ae_gen_kill}
Considerando un assegnamento $x := e$, dove $x$ è una variabile ed $e$
un'espressione.
\begin{itemize}
    \item \textbf{Killed}: Quando $x$ viene assegnata, tutte le espressioni
    che contengono $x$ non sono più valide (vengono "uccise") perché il
    valore di $x$ è cambiato.
    \[ kill(x := e) = \{ a \in AE \mid x \in FV(a) \} \]
    Dove $FV(a)$ è l'insieme delle variabili libere in $a$.
    
    \item \textbf{Gen}: Quando $x$ viene assegnata, l'espressione
    $e$ calcolata diventa disponibile, a condizione che $x$ non compaia
    nell'espressione stessa (altrimenti il valore appena calcolato verrebbe
    subito invalidato dall'assegnamento stesso).
    \[ gen(x := e) = \{ a \in AE \mid x \notin FV(a) \} \]
\end{itemize}
\end{definizione}

\paragraph{Equazioni di flusso}

Per ogni blocco $i \in nodes(G)$ del CFG, definiamo lo stato entrante
$entry(i)$ e lo stato uscente $exit(i)$.

\begin{definizione}{Stato entrante}{ae_entry}
Lo stato entrante di un blocco $i$ è dato dall'intersezione degli stati
uscenti di tutti i suoi predecessori:
\[
entry(i) = \bigcap \{exit(j) \mid (j, i) \in edges(G)\}
\]
la scelta dell'intersezione ($\cap$) è determinata dalla natura \textbf{must}
dell'analisi: un'espressione è disponibile solo se è disponibile lungo tutti
i cammini che raggiungono il blocco.

\end{definizione}

\begin {nota}{Blocco iniziale}{ae_entry_initial}
Se $i$ non ha predecessori (blocco iniziale), allora $entry(i) = \emptyset$.
\end{nota}

\begin{definizione}{Stato uscente}{ae_exit}
Lo stato uscente di un blocco $i$ è calcolato applicando la funzione di
trasferimento allo stato entrante:
\[
exit(i) = (entry(i) \setminus kill(stmt(G, i))) \cup gen(stmt(G, i))
\]

Ovvero:
\begin{enumerate}
    \item Si parte dalle espressioni disponibili in ingresso: $entry(i)$
    \item Si rimuovono le espressioni invalidate: $\setminus kill(stmt(G, i))$
    \item Si aggiungono le espressioni generate: $\cup gen(stmt(G, i))$
\end{enumerate}
\end{definizione}


Da queste funzioni è possibile riscrivere le funzioni $gen$ e $kill$ in modo
più generale:
\begin{itemize}
    \item $kill(x := \texttt{e}) = \{ a \in AE \mid x \in FV(\texttt{a}) \}$
    \item $gen(x := \texttt{e}) = \{ a \in AE \mid x \notin FV(\texttt{a}) \}$
\end{itemize}

\subsubsection{Altre analisi di dataflow}

Come accennato nella nota (\Cref{nt:types_of_df}), le analisi di dataflow
possono essere classificate secondo vari criteri.
In aggiunta all'analisi delle available expressions, esistono numerose altre
analisi di dataflow utilizzate per diversi scopi nell'ottimizzazione del
codice. Alcuni esempi di altre analisi viste nel corso di linguaggi interpreti e
compilatori di dataflow includono:
\begin{itemize}
    \item \textbf{Live variables}
    \item \textbf{Reaching definitions}
    \item \textbf{Very busy expressions}
\end{itemize}
L'immagine seguente classifica alcune di queste analisi in base alla direzione del
flusso e alla semantica della proprietà

\begin{figure}[H]
  \centering
  \includegraphics[width=0.8\textwidth]{images/th_04/02.png}
  \caption{L'immagine illustra la classificazione di alcune analisi di dataflow.}
  \label{fig:th_04_02}
\end{figure}

\begin{definizione}{Reaching definitions}{reaching_definitions}
``For each program point, which assignments may have been made and not
overwritten, when program execution reaches this point along some path''

Un'assegnazione $d: x := e$ al blocco $i$ raggiunge (reaches) un punto del
programma $p$ se esiste almeno un cammino dal blocco $i$ al punto $p$ lungo
il quale la variabile $x$ non viene ridefinita.

Quest'analisi è di tipo \textbf{forward} e \textbf{may}; l'informazione
si propaga in avanti nel CFG ed è sufficiente che un'assegnazione raggiunga
il punto lungo almeno un cammino.
\end{definizione}

\begin{definizione}{Live variables}{live_variables}
``For each program point, which variables may be live at the exit from the
program point. A variable is live if its content will be read in some path
starting from the program point''

Una variabile $x$ è viva (live) in un punto del programma $p$ se esiste
almeno un cammino che parte da $p$ lungo il quale il valore di $x$ viene
utilizzato prima di essere ridefinito.

Quest'analisi è di tipo \textbf{backward} e \textbf{may}; l'informazione
si propaga all'indietro nel CFG (dai successori ai predecessori) ed è
sufficiente che la variabile sia utilizzata lungo almeno un cammino
futuro.
\end{definizione}

\begin{definizione}{Very busy expressions}{very_busy_expressions}
``For each program point, which expressions must be very busy at the exit
from the point. An expression is very busy if it will be always used in all
paths starting from a label before any of the variables occurring in it are
redefined''

Un'espressione $e$ è molto impegnativa (very busy) in un punto del programma
$p$ se lungo ogni cammino che parte da $p$ l'espressione $e$ viene
calcolata prima che qualsiasi variabile che occorre in $e$ venga ridefinita.

Quest'analisi è di tipo \textbf{backward} e \textbf{must}; l'informazione
si propaga all'indietro nel CFG e l'espressione deve essere utilizzata lungo
tutti i cammini futuri.
\end{definizione}

\subsection{I reticoli nelle analisi di dataflow}

Tutte le analisi di dataflow discusse operano su dei reticoli completi
(\Cref{def:complete_lattice}).

\begin{definizione}{Reticolo delle analisi dataflow}{df_lattices}
Ogni analisi dataflow può essere definita come una tupla che rappresenta un
reticolo completo (\Cref{def:complete_lattice}).
Per le analisi viste, i reticoli sono:
\begin{enumerate}
    \item \textbf{Available expressions}: $\langle \mathcal{P}(AE), \supseteq, \emptyset, AE, \cap, \cup \rangle$.
    \item \textbf{Live variables}: $\langle \mathcal{P}(Var), \subseteq, \emptyset, Var, \cup, \cap \rangle$.
    \item \textbf{Reaching definitions}: $\langle \mathcal{P}(Def), \subseteq, \emptyset, Def, \cup, \cap \rangle$.
    \item \textbf{Very busy expressions}: $\langle \mathcal{P}(BE), \supseteq, \emptyset, BE, \cup, \cap \rangle$.
\end{enumerate}

Affinché l'algoritmo iterativo termini sicuramente, il reticolo deve soddisfare
la ascending chain condition (\Cref{def:acc}).
\end{definizione}

\begin{nota}{Modellazione analisi statiche}{df_modeling}
Non tutte le analisi statiche possono essere modellate come analisi di dataflow.
\end{nota}

\subsection{Analisi distributive}

Le analisi dataflow classiche appartengono alla categoria delle analisi
distributive
Questa proprietà matematica è fondamentale perché garantisce che l'analisi
possa essere risolta con algoritmi standard efficienti, ma ne limita anche la
potenza espressiva.

\begin{definizione}{Distributività}{distributivity}

Sia $\langle \mathcal{P}(X), \subseteq, \emptyset, X, \cup, \cap \rangle$ un
reticolo completo (\Cref{def:complete_lattice}) e $f: \mathcal{P}(X) \to \mathcal{P}(X)$
una funzione (\Cref{def:function}). $f$ si dice \textbf{distributiva} se:
\[ f(x \cup y) = f(x) \cup f(y) \]
per ogni $x, y \in \mathcal{P}(X)$.

\end{definizione}

\begin{nota}{Intuizione}{distributive_analysis}

Quanto descritto nella definizione sopra (\Cref{def:distributivity})
in termini intuitivi, significa che il risultato dell'analisi è identico sia
che si uniscano le informazioni prima di eseguire l'istruzione, sia
che le si uniscano dopo averla eseguita su ogni singolo ramo.

\end{nota}

\begin{nota}{La potenza delle analisi distributive}{power_of_distributive_analyses}

Le analisi distributive tengono traccia di \textbf{come} la computazione avviene
(es. variabili che sono assegnate ed espressioni che vengono calcolate)
all'interno del programma e non del \textbf{cosa} il programma calcola.

Per questo motivo le analisi complesse non sono generalmente distributive e molto
spesso non possono essere modellate come analisi di dataflow.

\end{nota}

\subsubsection{Constant propagation}

Un esempio di analisi non distributiva è la constant propagation.

\begin{definizione}{Constant propagation}{constant_propagation}
``
For each program point, whether or not a variable has a
constant value whenever an execution reaches that point
''

Un reticolo per la constant propagation può essere definito come:
\[
CP = \langle \mathbb{Z} \cup \{\bot\} \cup \{\top\}, \sqsubseteq, \bot, \top, \sqcup, \sqcap \rangle
\]

\begin{figure}[H]
  \centering
  \includegraphics[width=0.6\textwidth]{images/th_04/03.png}
  \caption{L'immagine illustra il reticolo per la constant propagation.}
  \label{fig:th_04_03}
\end{figure}

\end{definizione}
